\cotefile{22}
\scanpage{17}
\mspage{31}

En particulier, on obtient:

\uline{Corollaire} \deleted{3}\handwritten{4}.12. Soient $A\handwritten{\to}B$ un homomorphisme local d'anneaux locaux noethériens, tel\deleted{s} que $\underline{r}(B)=\underline{r}(A)B$ et \deleted{que l'extens} soit $K/k$ l'extension résiduelle. Conditions équivalentes: 1) $B$ est complet et formellement simple sur $A$. 2) $B$ est \deleted{\unsure{pl|p}}complet, plat sur $A$, et $K/k$ est séparable. \handwritten{3}) $B$ est une algèbre de Cohen sur $A$. De plus, une extension séparable $K$ de $k$ étant donnée, il existe une \deleted{\unsure{q|a}}algèbre de Cohen $B$ sur $A$ se réduisant suivant $K$, et $B$ est unique à isomorphisme de $A$-algèbres près.

La rédaction commence à faire un peu touffu, j'espère que tu pourras la rendre plus lisible. D'ailleurs, je me répens un peu de la définition \deleted{3}\handwritten{4}.7. telle quelle, il n'est pas vrai\deleted{\unsure{m|me}}ment naturel de supposer $B$ séparée et complète pour sa topologie propre, mais simplement pour sa topologie $IB$-adique (ce qui est moins fort); mais une telle modification demanderait des réajustement analogues pour \deleted{3}\handwritten{4}.2., \deleted{3}\handwritten{4}.6. et la suite. Je t'en laisse donc le choix, pensant que le choix adopté n'aura pas de grande importance pratique pour la suite : dans \deleted{3}\handwritten{4}.9. à \deleted{3}\handwritten{4}.12. cela obligera simplement à supprimer à certains endroits le mot \textquotedbl{}complet\textquotedbl{},\inserted{ou \deleted{\unsure{e|a}}\unsure{à|è|ă} le remplacer par un autre comple\unsure{t|\,}} \deleted{\deleted{en particulier \unsure{\textquotedbl{}Cohen\textquotedbl{}|Cohen} deviendra} synonyme dans le contexte des anneaux locaux noethériens de \deleted{\textquotedbl{}formellement} simple\textquotedbl{} , - sous réserve de l'embêtante hypothèse de \deleted{multiplicité} radicielle. \deleted{Il est vrai} Ce serait plutôt un avantage, \deleted{mais il est vrai par contre que cela ne s'harmoniserait plus} de plus les \deleted{anneaux de Cohen} sont}

\handwritten{\deleted{\uline{Corollaire \unsure{4.13|3.13}} $A\to B$ hom. loc. d'anneaux locaux noeth, munis de leurs topologies résiduelles. On suppose $k(B)/k(A)$ à mult. dép. finie. Alors conditions \unsure{équivalentes|équival.}}}
% archive sheet number: 16
